% !PRINTED 674
\weberpage{674}
%<<CONT
Quadrate der Determinante der Linearformen $\eta_s$, d.\,h. der Grundzahl des K\"orpers $\Omega$.

Daraus folgt, dass man eine von den $c_s$ unabh\"angige, nur durch den K\"orper $\Omega$ bestimmte Zahl $A$ finden kann, von der Art, dass, was auch die $c_s$ sein m\"ogen, $f(x_1,x_2,\dots,x_n)$ f\"ur ganzzahlige nicht verschwindende $x$ unter $A$ heruntersinkt. Da $f$ die Summe von positiven Summanden ist, so gilt dasselbe von jedem dieser Summanden, und damit ist der folgende Satz bewiesen:

3. \gesperrt{Ist $c_s$ ein beliebiges, den Bedingungen (5), (6) gen\"ugendes System reeller positiver Zahlen, so giebt es eine durch den K\"orper $\Omega$ allein bestimmte reelle Zahl $A$ von der Art, dass man immer eine ganze Zahl $\eta$ des K\"orpers $\Omega$ bestimmen kann, die mit ihren conjugirten Zahlen $\eta_s$ den Bedingungen gen\"ugt}
\weq{7}{|\eta_s| < c_s A.}
Beil\"aufig folgt hieraus, dass man in jedem algebraischen K\"orper $\Omega$, ausgenommen dem rationalen und dem imagin\"aren quadratischen, von Null verschiedene ganze Zahlen finden kann, deren absoluter Werth unter jede gegebene Grenze heruntersinkt.

Wir wollen den Ausdruck des Satzes 3. durch Einf\"uhrung der conjugirten Logarithmen von $\eta$ etwas umformen. Wir setzen $\log A = k$, so dass $k$ eine durch $\Omega$ bestimmte reelle Zahl ist, ferner
\weq{8}{c_s = e^{\frac{\gamma_s u}{\delta_s}},}
worin $u$ eine beliebige reelle Gr\"osse sein kann, und die Zahlen $\gamma_s$ wegen (5) und (6) der Bedingung gen\"ugen:
\weq{9}{\gamma_1 + \gamma_2 + \cdots + \gamma_\nu = \sum_{1,\nu}^{i} \gamma_i = 0.}
Da \"uberdies
\[
|\eta_s| = e^{\frac{\lambda_s}{\delta_s}}
\]
ist, so folgt aus (7) und (8):
\weq{10}{\lambda_s - \gamma'_s u < k\delta_s.}
% [sic] The printed prime has no referent; the unprimed \gamma_s is required by (8), (9), and the following sum.
Die Summe der $\nu$ Ausdr\"ucke auf der linken Seite dieser Ungleichung ist wegen (9) gleich $\Sigma\lambda_s$, also nach dem Satze 1. nicht negativ, und wir finden
\[
0 \leqq \sum^{s} (\lambda_s - \gamma_s u) < k n.
\]
